\chapter{Introduction}
\label{sec:intro}

\section{Motivation}
\label{sec:intro:mo}

Business process is valuable asset of an organization. But most of them are store in unstructured format, e.g., Word documents, PDF files. Munal translate from unstructured documents to structured process models is time-consuming and error-prone.
In the era of Gen-AI, LLM has shown as a promising solution to support creating proces modeling efficiently. However, it still cannot fully replace human in process modeling. Hybrid human-LLM process modeling is a promising  to leverage the strength of both human and LLM, and mitigate their weaknesses. 

However, existing tools and approaches still have limitations in terms of usability, integration, and functionality, especially in the use case for creating process models from documents. For example, user need to manually extract relevant information from documents and input it into the conversational chatbot. Besides, the separate interface of the document and the generated model make it hard to verify the result.

To address these limitations, this work proposes an integrated tool that supports process modeling from documents with a more seamless, efficient, and user-friendly workflow by offering the following core features: (1) allowing users to upload documents and select passages from the documents to serve as input for LLM-based process model generation directly, (2) providing a clear visual link between the text and the generated model through display document and model side by side, and (3) allowing users to establish process-subprocess relationships between models across the documents.

\section{Research Questions}
\label{sec:intro:rq}

This work addresses the following research questions:

RQ1 How do existing approaches and tools support process modeling from documents?

This question investigates how existing approaches and tools support process modeling from documents by identifying their functionalities and capabilities. It helps guide the design of the proposed system.

RQ2 How can an integrated system be designed to support process modeling from documents?

This question addresses the design of the proposed system. It focuses on what functionalities should be integrated and how to integrate them, aiming to enable a seamless modeling workflow with high usability.

RQ3 How do users perceive the proposed system?

This question evaluates the proposed system from the users’ perspective. It aims at understanding how users perceive the system’s benefits and limitations. Its results provide insights into the system’s strengths and weaknesses and help identify potential design issues and areas for improvement.


\section{Contribution}
\label{sec:intro:con}
1. The proposed and developed tool as an artefact of this work.
2. A proof-of-concept that demonstrates the feasibility of the proposed tool and its potential to support process modeling from documents through the improved modeling workflow.
\section{Methodology}
\label{sec:intro:meth}

This work follows the three interconnected cycles of DSR\cite{hevner2007} proposed by Hevner:
The relevance cycle: the first cycle is done through the communication with modeling experts, who proposes the tool and eliciate the requirements in section \ref{sec:solution:req} for the tool. 
The rigor cycle: it is done throgh the review in section \ref{sec:rel}.
The design cycle: the first cycle is done through the delivered solution of RQ2 and evaluation result of RQ3. Through the analysis and discussion of the evalution result, we can identify the strengths and weaknesses of the proposed tool, and further improve the design of the tool in the next iteration.

\section{Evaluation}
\label{sec:intro:ev}

The evaluation of the proposed system is done with 4 experts in qualitative way. It consisted of task-based activities and a follow-up questionnaire to collect user feedback on the tool's usability and perceived usefulness.

\section{Structure}
\label{sec:intro:struct}

The remainder of this thesis is structured as follows: Chapter~\ref{sec:rel} reviews the related work on process modeling from documents and LLM-based process modeling. Chapter~\ref{sec:solution} presents the design of the proposed system, including its requirements, architecture, and functionalities. Chapter~\ref{sec:eval} describes the evaluation of the proposed system, including the methodology, results. Chapter~\ref{sec:discussion} discusses the implications of the findings and potential future work. Finally, Chapter~\ref{sec:conclusion} concludes the thesis and outlines future work.
